\section{Related Work}
\label{sec:related}

Understanding why deduplication's benchmark effects are contamination-proportional rather
than uniform requires drawing on three lines of prior research: (1) the documented effects
of deduplication on language model performance, (2) benchmark contamination detection and
measurement, and (3) the Pythia controlled training infrastructure.

\subsection{Deduplication and Training Data Quality}

Data deduplication has been established as a beneficial preprocessing step for language
model training. \citet{Lee2022Deduplicating} provided the first systematic study showing that
deduplicating training data using MinHash and exact substring methods reduces memorization
and improves average GPT-2-scale model performance across multiple benchmarks.
This foundational result motivated widespread adoption of deduplication in subsequent
corpus construction efforts including C4 \citep{Raffel2020Exploring}, RefinedWeb, Dolma
\citep{Soldaini2024Dolma}, and DCLM.

However, \citet{Lee2022Deduplicating} reported aggregate averages across benchmarks, and did not
test whether the improvement was uniform across all benchmarks or concentrated on
specific benchmarks with higher contamination. \citet{Muennighoff2023Scaling} analyzed
the impact of data repetition on model performance, finding diminishing returns at
high repetition rates, but focused on training dynamics rather than contamination-proportional
benchmark effects. Our work is the first to test the contamination-correlation prediction
directly: does the per-benchmark improvement from deduplication correlate with each
benchmark's contamination level in the original corpus?

\citet{Biderman2023Pythia} introduced the Pythia model suite with an explicit Pile vs.\
dedup-Pile controlled comparison, reporting benchmark results that showed mixed directions
across benchmarks. That paper did not perform contamination-accuracy correlation analysis
and used step-matched comparisons that we show introduce a volume confound
(our H-M4 result: $\Delta r = 0.093$ relative to token-count matching). Our work extends
\citet{Biderman2023Pythia} by providing the mechanistic analysis that was absent from the
original study and correcting the methodological comparison.

\subsection{Benchmark Contamination Detection}

The problem of benchmark contamination has received increasing attention as LLMs have been
shown to achieve high benchmark scores partly through training data memorization rather
than genuine task understanding.

\citet{Brown2020Language} first characterized contamination in GPT-3 training data using n-gram
overlap analysis, acknowledging that benchmark test sets may appear in web-crawled
training corpora. The GPT-4 Technical Report adopted a 13-gram overlap methodology
for characterizing contamination in training data, establishing a standard that
subsequent work has followed.

\citet{Shi2023Detecting} introduced Min-$k$\% Probability (min-$k$\%), a likelihood-ratio-based
method for detecting whether a given text was present in an LLM's pretraining data.
By computing the minimum token probabilities over $k$\% of tokens in a sequence,
min-$k$\% provides a membership inference signal for pretraining data detection. \citet{Shi2023Detecting}
validated this method on multiple benchmarks, showing it outperforms perplexity-based
detection. Our experiment applies min-$k$\% to the specific Pile/dedup-Pile distinction
and finds an unexpected direction reversal at Pythia-1B scale, raising questions about
min-$k$\%'s sensitivity for corpus variants that are partially overlapping rather than
fully seen vs.\ unseen.

\citet{Carlini2021Extracting} demonstrated that memorization in language models scales with model
size: larger models memorize more training examples verbatim. This scaling relationship
provides the theoretical basis for our observation that the min-$k$\% memorization signal
may require models larger than 1B parameters to manifest in the predicted direction for the
Pile/dedup-Pile distinction.

\citet{Golchin2023Time} proposed data contamination quiz methods for detection in
instruction-tuned models; this line of work focuses on detecting contamination post-hoc
rather than characterizing its effect on benchmark performance as a function of contamination
level, which is our primary focus.

\subsection{The Pile, dedup-Pile, and Pythia}

The Pile \citep{Gao2020Pile} is an 825GB English text corpus constructed by EleutherAI
from 22 diverse data sources. Its deduplication variant, dedup-Pile, applies exact
substring deduplication to remove repeated content, reducing the corpus by approximately
15\% in token count. Both corpora were used to train the Pythia model suite
\citep{Biderman2023Pythia} under otherwise identical conditions, creating a unique
natural experiment for isolating deduplication's causal effects.

The Pythia suite provides 154 intermediate checkpoints per model size (70M--12B
parameters), enabling the token-count matching methodology that our study depends on.
Without this checkpoint granularity, identifying Pile checkpoints at the same token count
as the dedup-Pile final checkpoint would not be feasible.

The OLMo suite \citep{Groeneveld2024OLMo} and Dolma corpus \citep{Soldaini2024Dolma} provide
an alternative open-weights model family with documented curation choices. However, OLMo
and Pythia use different architectures (OLMo vs.\ GPT-NeoX), preventing a clean causal
interpretation of any cross-family comparison. We focus exclusively on the within-family
Pythia controlled comparison and note cross-family extension as future work.

\subsection{Scaling Laws and Token-Volume Effects}

\citet{Hoffmann2022Training} (Chinchilla) established that optimal model training requires
matching training token count to model parameters via compute-optimal scaling laws.
Their analysis implies that models trained on the same step count but different token
volumes (due to corpus size differences) will have different effective training data
amounts. This provides the theoretical grounding for our token-count matching
methodology: Pile and dedup-Pile models trained to the same step count have seen
different numbers of tokens, introducing a volume confound that step-matching fails to correct for.

\subsection{Our Position}

Our work sits at the intersection of these three research streams. Unlike \citet{Lee2022Deduplicating},
which reports aggregate deduplication benefits, we decompose the benchmark-level
effects and test a contamination-proportionality prediction. Unlike \citet{Shi2023Detecting}
and \citet{Carlini2021Extracting}, which focus on memorization detection, we use contamination
estimates as a predictor variable to explain benchmark performance changes under a
curation intervention. And unlike \citet{Biderman2023Pythia}, which reports Pythia
benchmark results without mechanistic analysis or volume-confound correction, we
provide both the contamination-accuracy correlation and the token-count matching
methodology that recovers the full contamination signal.
